\newpage

% =========================================================================
% ACESSO AO CÓDIGO-FONTE E REPOSITÓRIO GITHUB
% =========================================================================
\apendice{Acesso ao Código-Fonte e Repositório GitHub}
\label{ap:github}

Para garantir a reprodutibilidade dos experimentos e disponibilizar o código-fonte desenvolvido, todo o material está organizado em repositório público no GitHub.

O repositório \textbf{Network Flow} está estruturado nos seguintes diretórios principais:
\begin{itemize}[leftmargin=1.5em, itemsep=0.2em]
\raggedright
    \item \textbf{Fluxo Máximo} (\path{implementations/flow-network/}): Interface base \lstinline{FlowNetwork} e algoritmos concretos (\lstinline{FordFulkerson}, \lstinline{EdmondsKarp}, \lstinline{Dinic}, \lstinline{MPM}, \lstinline{PushRelabel}, \lstinline{PushRelabelImproved}, \lstinline{Kuhn} e \lstinline{HopcroftKarp});
    \item \textbf{Fluxo de Custo Mínimo} (\path{implementations/cost-network/}): Interface base \lstinline{CostNetwork} e algoritmos concretos (\lstinline{CycleCanceling}, \lstinline{MinimumMeanCycleCanceling}, \lstinline{CancelAndTighten}, \lstinline{SuccessiveShortest}, \lstinline{SuccessiveShortestDijkstra}, \lstinline{NetworkSimplex} e \lstinline{KuhnMunkres});
    \item \textbf{Problemas} (\path{implementations/problems/}): Resoluções para problemas combinatórios clássicos e de maratona (como \textit{Minimum Vertex Cover}, emparelhamentos, caminhos disjuntos, fechamento de peso máximo e demandas mínimas);
    \item \textbf{Aplicações Práticas} (\path{applications/image-segmentation/}): Segmentação de imagens via corte mínimo em grafos (\textit{Graph Cuts});
    \item \textbf{Experimentos} (\path{experiments/}): Scripts de automação de benchmarks, leitura de instâncias DIMACS e geração de tabelas e gráficos experimentais.
\end{itemize}

O código-fonte completo pode ser consultado, clonado e executado através do endereço:

\vspace{0.4cm}
\begin{center}
    \fbox{\url{https://github.com/GabrielFrigo4/network-flow}}
\end{center}
\vspace{0.4cm}

O repositório contém o histórico de desenvolvimento e as instruções para compilação e reprodução dos experimentos.
